% !PRINTED 360
\weberpage{360}
Nun ist die nothwendige und hinreichende Bedingung dafür, dass eine ternäre quadratische Form in zwei lineare Factoren zerfalle, die, dass die Determinante der quadratischen Form verschwinde (Bd. I, \S.\,56).

Drücken wir die Formen $u, v$ durch $x_1, y_1$ und irgend eine dritte Variable $y$ aus, und bezeichnen die Co\"efficienten in $v$ und $u$ mit $a_{ik}, b_{ik}$, so dass $2a_{23}, 2b_{23}$ die Co\"efficienten von $x_1 y_1$ werden, so erhalten wir die Bedingung für das Zerfallen von $v + 2\lambda u + \lambda^2 x_1 y_1$ in der Form:
\weq{6}{\begin{vmatrix} a_{11}+2\lambda b_{11}, & a_{12}+2\lambda b_{12}, & a_{13}+2\lambda b_{13} \\ a_{21}+2\lambda b_{21}, & a_{22}+2\lambda b_{22}, & a_{23}+2\lambda b_{23}+\tfrac{1}{2}\lambda^2 \\ a_{31}+2\lambda b_{31}, & a_{32}+2\lambda b_{32}+\tfrac{1}{2}\lambda^2, & a_{33}+2\lambda b_{33} \end{vmatrix} = 0,}
und dies ist eine Gleichung $5^{\mathrm{ten}}$ Grades in Bezug auf $\lambda$, die uns also fünf Zerlegungen giebt:
\[
x_2 y_2,\ x_3 y_3,\ x_4 y_4,\ x_5 y_5,\ x_6 y_6.
\]
Ueber die Co\"efficienten in $f$ lässt sich, wie (6) zeigt, so verfügen, dass alle die so bestimmten Functionen $x_i y_i$ von einander verschieden sind, und daraus folgt, wie oben gezeigt, dass sie verschieden bleiben, so lange die Curve $f$ keine singulären Punkte hat. Es gilt also der folgende Satz:

1. \gesperrt{Zu jedem Paare von Doppeltangenten $x_1 y_1$ gehören fünf weitere Paare $x_i y_i$ von der Art, dass die acht Berührungspunkte von $x_1 y_1 x_i y_i$ auf einem Kegelschnitte liegen.}

Die sechs Paare von Doppeltangenten, die auf diese Weise bestimmt sind:
\[
x_1 y_1,\ x_2 y_2,\ x_3 y_3,\ x_4 y_4,\ x_5 y_5,\ x_6 y_6,
\]
wollen wir einen \gesperrt{Steiner'schen Complex} nennen\footnote{Es ist dafür bisher gewöhnlich der Ausdruck \gesperrt{Steiner'sche Gruppe} gebraucht. Da aber das Wort \glqq Gruppe\grqq{} in der Algebra eine ganz bestimmte andere Bedeutung hat, so ziehen wir es vor, diesen Ausdruck hier zu vermeiden.}.

Betrachten wir die Paare $x_1 y_1, x_2 y_2, x_3 y_3$ eines solchen Complexes, und setzen
\[
f = x_1 y_1 x_2 y_2 - u_1^2 = x_1 y_1 x_3 y_3 - u_2^2,
\]
%CONT>>
